% Generated draft: requires mathematical and source review.
\documentclass[11pt]{article}
\usepackage[margin=25mm]{geometry}
\usepackage{fontspec}
\setmainfont{Latin Modern Roman}
\usepackage{amsmath,amssymb,mathtools,unicode-math}
\setmathfont{Latin Modern Math}
\usepackage{enumitem,needspace,xurl,hyperref}
\setlength{\parindent}{0pt}
\setlength{\parskip}{5pt}
\setlength{\emergencystretch}{3em}
\newcommand{\N}{\mathbb{N}}
\newcommand{\Z}{\mathbb{Z}}
\newcommand{\Q}{\mathbb{Q}}
\newcommand{\R}{\mathbb{R}}
\newcommand{\C}{\mathbb{C}}
\newcommand{\F}{\mathbb{F}}
\newcommand{\A}{\mathbb{A}}
\newcommand{\PP}{\mathbb{P}}
\newcommand{\E}{\mathbb{E}}
\newcommand{\eps}{\varepsilon}
\newcommand{\dd}{\,\mathrm{d}}
\begin{document}
\section*{UnsolvedMath ID: 30000926}
\textbf{Source problem code:} catalogue entry
\subsection{Definitions and mathematical statement}
Let $A(\mathbf x,Q)$ be a cluster algebra with initial cluster $\mathbf x=(x_1,\ldots,x_r)$ and an initial quiver $Q$ that is not acyclic, meaning it has an oriented cycle. The source asks how to interpret the denominator of an arbitrary cluster variable relative to this initial seed, and whether every cluster variable admits the report's notion of a $T$-denominator.

For a Laurent expression $v=\sum_a c_a\mathbf x^a$, a standard signed denominator-vector convention is
\[
d_i(v)=-\min\{a_i:c_a\ne0\},\qquad
v=\frac{P(\mathbf x)}{\prod_i x_i^{d_i(v)}},
\]
where $P$ is not divisible by any $x_i$ and the coefficient system is fixed. This convention gives $d(x_j)=-e_j$ for an initial variable. A convention requiring a nonnegative denominator monomial treats such variables differently; the record does not specify which convention or initial-variable exception occurs in its definition.

The requested categorical formula is
\[
d_y(v)=\dim\operatorname{Hom}_{\mathcal C_Q}(T_y,M_v),
\qquad T=\bigoplus_yT_y,
\]
where $T$ is the corresponding cluster-tilting object and $M_v$ is the object corresponding to $v$. More explicitly, the question asks whether all cluster variables have a denominator with these exponents in the sense of the report's $T$-denominator definition. It also retains the broader request for a denominator interpretation when the initial seed is nonacyclic.

The supplied record does not define $\mathcal C_Q$, the correspondence $v\mapsto M_v$, the base field for the Hom dimension, or the precise $T$-denominator definition. A conventional cluster category built from an acyclic quiver cannot simply be assigned to an arbitrary cyclic quiver without further construction or hypotheses. Nonacyclicity of an initial seed also does not specify whether the cluster algebra is mutation-equivalent to an acyclic one. The full range of algebras and the exceptional treatment of initial variables are therefore unresolved: nonnegative Hom dimensions cannot directly equal the negative initial-variable vectors in the signed convention above. No shifts, corrections or special-case formula are invented to remove this gap.

The supplied assessment checked 2026-08-21 labels the problem partially solved. It reports denominator formulas in finite, affine, and several surface or categorical settings, but no theorem settling the broad nonacyclic initial-seed question for all cluster algebras. This is dated attributed triage, not an independently verified universal status statement.
\subsection{Short English statement}
Interpret denominator vectors relative to a nonacyclic initial seed and ask whether the report's $T$-denominator formula using Hom dimensions holds for every cluster variable. The category, correspondence and initial-variable conventions needed to state that claim precisely are missing.
\subsection{Sources}
\begin{itemize}
\item Local UnsolvedMath registry, permanent ID 30000926, source code catalogue entry; supplied complete source record, including its reference to the report's $T$-denominator definition and corrected-verified source-check note.
\item Representation Theory of Finite Dimensional Algebras (2008), Oberwolfach Reports, DOI 10.4171/owr/2008/09, pp. 401--472: \url{https://doi.org/10.4171/owr/2008/09}.
\item Affine denominator theorem: \url{https://doi.org/10.1016/j.jalgebra.2010.02.012}.
\item Finite-type denominator vectors: \url{https://arxiv.org/abs/1302.1052}.
\end{itemize}
\end{document}
